\documentclass[11pt]{article}
\usepackage[utf8]{inputenc}
\usepackage{amsmath,amssymb,amsthm}
\usepackage[margin=1in]{geometry}
\usepackage{enumitem}
\usepackage{hyperref}

\begin{document}
% ==================== PAGE 13 ====================
\noindent \textbf{Theorem:} If $\vert{}G\vert{} = p^n$, where $p$ is a prime then $\vert{}Z(G)\vert{} \ge p$.

\medskip
\noindent \textbf{Proof:} \\
Let $\vert{}Z(G)\vert{} = r$. \\
By class equation:
\[\vert{}G\vert{} = \vert{}Z(G)\vert{} + \sum_{a \notin Z(G)} \left\vert{}\frac{G}{N(a)}\right\vert{} \tag{A}\]
$a \notin Z(G) \implies N(a) \ne G$ \\
$\implies N(a)$ is a proper subgroup of $G$. \\
By Lagrange's theorem, $\vert{}N(a)\vert{} \mid \vert{}G\vert{}$
\[\Rightarrow \vert{}N(a)\vert{} = p^{\alpha_a}, \quad \alpha_a < n\]
\[\left\vert{}\frac{G}{N(a)}\right\vert{} = \frac{\vert{}G\vert{}}{\vert{}N(a)\vert{}} = \frac{p^n}{p^{\alpha_a}} = p^{n - \alpha_a}, \quad n - \alpha_a > 0\]
\[\Rightarrow p  \mid  \left\vert{}\frac{G}{N(a)}\right\vert{}\]
\[\Rightarrow p  \mid  \sum_{a \notin Z(G)} \left\vert{}\frac{G}{N(a)}\right\vert{}\]
Also, $p \mid \vert{}G\vert{}$ (as $\vert{}G\vert{} = p^n$)
\[\Rightarrow p  \mid  \left(\vert{}G\vert{} - \sum_{a \notin Z(G)} \left\vert{}\frac{G}{N(a)}\right\vert{}\right)\]
From (A):
\[\Rightarrow p \mid \vert{}Z(G)\vert{}\]
Since $e \in Z(G) \implies \vert{}Z(G)\vert{} \ge 1$. \\
$\therefore p \le \vert{}Z(G)\vert{} \implies \vert{}Z(G)\vert{} \ge p$.

\bigskip
\noindent $\text{Cl}((12)) = \{(12), (13), (14), (23), (24), (34)\}$ \\
$\text{Cl}((123))$

\newpage
% ==================== PAGE 14 ====================
\noindent \textbf{Theorem:} A group of order $p^2$, $p$ is prime, is abelian.

\medskip
\noindent \textbf{Proof:} \\
Let $\vert{}G\vert{} = p^2$. \\
$\Rightarrow \vert{}Z(G)\vert{} \ge p$. \\
So, $\vert{}Z(G)\vert{} = p$ or $p^2$.

\noindent If $\vert{}Z(G)\vert{} = p$, then
\[\left\vert{}\frac{G}{Z(G)}\right\vert{} = \frac{p^2}{p} = p\]
$\Rightarrow G/Z(G)$ is cyclic \\
$\Rightarrow G$ is abelian \\
$\Rightarrow G = Z(G) \implies \vert{}G\vert{} = \vert{}Z(G)\vert{} = p^2$. \\
$\therefore \vert{}Z(G)\vert{} \ne p$ (contradiction).

\noindent So, $\vert{}Z(G)\vert{} = p^2 = \vert{}G\vert{}$ \\
$\Rightarrow G = Z(G)$ \\
$\therefore G$ is abelian.

\bigskip
\noindent \textbf{Note:} If $H \le G$ s.t. $H \le Z(G)$ then $H \trianglelefteq G$.

\bigskip
\noindent \textbf{Theorem:} Let $\sigma$ be a permutation in $S_n$ and let $m_1, m_2, \dots, m_r$ be the distinct integers which appear in the cycle type of $\sigma$ (including 1-cycles). Assume that $\sigma$ has $k_i$ cycles of length $m_i$ for each $i = 1, 2, \dots, r$ so that $\sum_{i=1}^r k_i m_i = n$. \\
Then number of conjugates of $\sigma$ is
\[\frac{n!}{(k_1! \, m_1^{k_1})(k_2! \, m_2^{k_2}) \cdots (k_r! \, m_r^{k_r})}\]

\newpage
% ==================== PAGE 15 ====================
\noindent \textbf{Example:} In $S_4$ \\
No. of cycles of 2 length in $S_4$:
\[\frac{4!}{(1! \, 2^1)(2! \, 1^2)} = \frac{24}{2 \times 2} = 6\]
No. of cycles of length 3:
\[\frac{4!}{(1! \, 3^1)(1! \, 1^1)} = \frac{24}{3} = 8\]
In $S_5$:
\[5! = 5 \times 4 \times 3 \times 2 \times 1 = 120\]
\[(3! \, 1^3)(1! \, 2^1) = 6 \times 2 = 12 \implies \frac{120}{12} = 10\]

\bigskip
\noindent \textbf{Ring:} A system $(R, +, \cdot)$ consisting of a non-empty set $R$ and two internal binary operations `$+$' and `$\cdot$' in $R$ is called a \textbf{ring} if:
\begin{enumerate}
    \item $(R, +)$ is an abelian group
    \item $(R, \cdot)$ is a semi-group
    \item Distributive laws hold:
    \begin{align*}
    a \cdot (b + c) &= a \cdot b + a \cdot c \\
    (a + b) \cdot c &= a \cdot c + b \cdot c \quad \forall\, a, b, c \in R
    \end{align*}
\end{enumerate}

\noindent If $(R, \cdot)$ has identity $\implies$ \textbf{Ring with unity}. \\
If $(R, \cdot)$ is abelian $\implies$ \textbf{Commutative ring}.

\bigskip
\noindent \textbf{Note:}
\begin{enumerate}
    \item `$+$' is usually called addition and `$\cdot$' is usually called multiplication.
    \item Identity of $(R, +)$ is called the \textbf{zero element} of ring and is denoted as $0$.
    \item The identity w.r.t. `$\cdot$' is known as \textbf{unity} (identity element) of the ring and is denoted as $1$.
\end{enumerate}

\newpage
% ==================== PAGE 16 ====================
\begin{enumerate}[resume]
    \item If `$\cdot$' is commutative then $R$ is called as \textbf{commutative ring}.
    \item If the identity exists in $R$ then $R$ is called as \textbf{ring with identity} or \textbf{ring with unity}.
\end{enumerate}

\bigskip
\noindent \textbf{Example:}
\begin{enumerate}
    \item $(\mathbb{Z}, +, \cdot)$, $(\mathbb{Q}, +, \cdot)$, $(\mathbb{R}, +, \cdot)$, $(\mathbb{C}, +, \cdot)$ all are commutative ring with unity.
    \item $(\mathbb{Z}_n, +, \cdot)$ is a commutative ring with unity.
    \item $\mathbb{Z}[\sqrt{p}] = \{a + b\sqrt{p} \mid a, b \in \mathbb{Z}, \; p \text{ is prime}\}$. Then it is a commutative ring with unity.
    \item $\mathbb{Z}[i] = \{a + bi \mid a, b \in \mathbb{Z}\}$ is also a commutative ring with unity (\textbf{Ring of Gaussian Integers}).
    \item $M_n(\mathbb{R}) = \text{set of all } n \times n \text{ matrices with real entries}$. \\
    $(M_n(\mathbb{R}), +)$ is group. \\
    $(M_n(\mathbb{R}), \cdot)$ is a semi-group. \\
    $A \cdot (B + C) = AB + AC$ \\
    $(A + B) \cdot C = AC + BC$ \\
    $\therefore (M_n(\mathbb{R}), +, \cdot)$ is a non-commutative ring with unity.
\end{enumerate}
\end{document}
